\slajdid{02-s026}
\begin{frame}[label=02-s026]
\gusttekst
\setlength{\abovedisplayskip}{3pt}\setlength{\belowdisplayskip}{3pt}
Za realizaciju funkcije u obliku proizvoda zbirova
\[
\begin{aligned}
F&=(C+B+A)(C+B+\bar A)(C+\bar B+A)(\bar C+B+A)(\bar C+\bar B+\bar A)\\
&=\overline{\overline{(C+B+A)(C+B+\bar A)(C+\bar B+A)(\bar C+B+A)(\bar C+\bar B+\bar A)}}\\
&=\overline{\overline{(C+B+A)}+\overline{(C+B+\bar A)}+\overline{(C+\bar B+A)}+\overline{(\bar C+B+A)}+\overline{(\bar C+\bar B+\bar A)}}
\end{aligned}
\]
ili preko električnih, logičkih, šema

\centering
\input{slike/tikz/s026_nili_i_invertovani_ulazi.tex}
\(\longrightarrow\)
\begin{tikzpicture}[font=\fontsize{8}{9.5}\selectfont,x=6mm,y=4mm,baseline=(current bounding box.center),
 DEinvertor/.style={not gate US,draw,scale=.6,minimum width=5mm,minimum height=4mm},
 DEprvinivo/.style={nor gate US,draw,scale=.6,logic gate inputs=nnn,rotate=-90,minimum width=5mm,minimum height=5mm},
 DEdrug nivo/.style={nor gate US,draw,scale=.6,logic gate inputs=nnnnn,rotate=-90,minimum width=6mm,minimum height=7mm}]
\def\DEkrajvodova{8.7}
% Gornji vod para je prava vrednost; donji je komplementna.
\foreach \y/\low/\lab in {6.5/5.9/C,5.3/4.7/B,4.1/3.5/A}{
 \node[DEinvertor] (n\lab) at (.65,\y) {};
 \node[DEinvertor] (m\lab) at (2.0,\y) {};
 \draw[DEvod] (0,\y) node[left] {\(\lab\)}--(n\lab.input);
 \draw[DEvod] (n\lab.output)--(m\lab.input);
 \draw[DEvod] (m\lab.output)--(\DEkrajvodova,\y);
 \coordinate (t\lab) at ($(n\lab.output)!.5!(m\lab.input)$);
 \draw[DEvod] (t\lab)--(t\lab |- 0,\low)--(\DEkrajvodova,\low);
}

\foreach \i/\x in {1/3.1,2/4.4,3/5.7,4/7,5/8.3}{
 \node[DEprvinivo] (g\i) at (\x,2.1) {};
}
\foreach \i/\j/\y in {1/3/6.5,1/2/5.3,1/1/4.1,2/3/6.5,2/2/5.3,2/1/3.5,3/3/6.5,3/2/4.7,3/1/4.1,4/3/5.9,4/2/5.3,4/1/4.1,5/3/5.9,5/2/4.7,5/1/3.5}{
 \draw[DEvod] (g\i.input \j)--(g\i.input \j |- 0,\y);
}
\node[DEdrug nivo] (z) at (5.7,-.3) {};
\draw[DEvod] (g1.output)--(g1.output |- 0,.65)-|(z.input 5);
\draw[DEvod] (g2.output)--(g2.output |- 0,1.05)-|(z.input 4);
\draw[DEvod] (g3.output)--(z.input 3);
\draw[DEvod] (g4.output)--(g4.output |- 0,1.05)-|(z.input 2);
\draw[DEvod] (g5.output)--(g5.output |- 0,.65)-|(z.input 1);
\draw[DEvod] (z.output)--++(0,-.3) node[below] {\(F\)};

\end{tikzpicture}
\par
\raggedright
Znači moguće je bilo koju kombinacionu mrežu realizovati, izuzimajući invertore samo pomoću NILI logičkih kola. Ili generalnije ako invertor smatramo NILI logičkim kolom sa jednim ulazom: Bilo koju kombinacionu mrežu možemo realizovati korišćenjem samo NILI logičkih kola.
\end{frame}
